\begin{textsample}{POS Dim 3 – pa – Score 9.00 – pa/pa\_inf\_11.txt}  \label{ex:f3_pos_278}
3.1.5 A Relação Família e Escola A Educação \textbf{Infantil} enquanto primeira etapa da educação básica apresenta uma variação dos demais níveis de ensino quando apresenta a família como importante articuladora no processo de aprendizagem das \textbf{crianças}, pois é nesse primeiro momento de suas vidas que elas são afastadas física e temporariamente de suas relações parentescas.
Afirma -se que esse é o primeiro momento de transição que a \textbf{criança} passa ao “ deixar”[20 ] seu lar para começar a construir seu vínculo com o espaço de educação formal que chamamos de escola, no caso, de \textbf{creche} ou unidades de Educação \textbf{Infantil}.
Essa separação não pode em hipótese nenhuma transformar -se em rompimento, visto que a educação dela deve ser complementar e articulada com a que é oferecida pelas suas famílias, como destaca a LDB no Art. 29.
A Educação \textbf{Infantil}, primeira etapa da educação básica, tem como finalidade o desenvolvimento integral da \textbf{criança} até seis anos de idade, em seus aspectos físico, psicológico, intelectual e social, completando a ação da família e da comunidade ( BRASIL, 1996, n.p. ).
À escola cabe o papel de incluir a família no processo de desenvolvimento da \textbf{criança} frente ao projeto educacional, estimulando e contribuindo assim para que se estabeleça uma relação de confiança entre os sujeitos e entidades responsáveis pela educação da \textbf{criança} ; situação esta que possibilita ainda que a família participe efetivamente do que se propõe para os \textbf{cuidados} e o processo de aprendizagem na Educação \textbf{Infantil} e das ações que constam na proposta pedagógica dos espaços educacionais.
A Base propõe direitos de aprendizagem de conviver, \textbf{brincar}, participar, explorar, comunicar, conhecer -se, e esses direitos para serem consolidados devem buscar se pautar nas experiências de aprendizagem, experiências concretas na vida cotidiana que levam à aprendizagem da cultura, pelo convívio no espaço coletivo, e a produção de narrativas, individuais e coletivas, por meio de diferentes linguagens ( BRASIL, 2017a ).
A relação entre escola e família se mostra fundamental para favorecer a articulação entre as experiências ora citadas e os saberes que serão aprendidos para garantir o pleno desenvolvimento da \textbf{criança}, ou seja, o processo de aprendizado se constituirá dentro das interações sociais, visto que, segundo Vygotsky ( 1991 ) a \textbf{criança} é um sujeito histórico que se constitui a partir de interações e das relações sociais que são estabelecidas desde muito cedo por ela.
Cruz ( 2016 ) nos chama a tentar entender as expectativas tanto da escola quanto da família frente ao processo de aprendizado das \textbf{crianças}, particularmente ao que tange o processo de leitura e escrita, pois para a autora o entendimento sob as expectativas é fundamental “ [... ] para analisar as convergências e divergências destas com o que é proposto às \textbf{crianças} na Educação \textbf{Infantil} e estabelecer um diálogo mais profícuo em relação a esse tema ” ( CRUZ, 2016, p. 16 ).
Nesse contexto de relação exitosa em prol do desenvolvimento biológico, afetivo, emocional, cultural e social da \textbf{criança}, a escola precisa \textbf{demonstrar} solidez em seus conceitos e concepções acerca do que defende para a Educação \textbf{Infantil} e isso deve se refletir nos profissionais que cuidam e educam as \textbf{crianças}, para que assim se legitime, junto à família, que o \textbf{brincar}, o imaginário, a fantasia, os \textbf{desejos}, os pensamentos, as falas e os movimentos corporais são importantes e imprescindíveis para processo de aprendizagem.
O diálogo entre família e escola, historicamente se mostrou difícil, pois põe em evidência o julgamento por parte de cada um desses sujeitos faz um do outro, posto que as reuniões não facilitem a “ [... ] oportunidade de maior conhecimento e troca entre os profissionais da escola e os familiares das \textbf{crianças} ” ( CRUZ, 2016, p. 22 ).
Esse mesmo contexto de tensões e conflitos que se interpõe no cotidiano dos espaços educacionais e envolvem professores e famílias, retrata momentos de cooperação e parceria entre os mesmos sujeitos que buscam garantir o desenvolvimento integral da \textbf{criança}.
Para a constituição de uma legítima parceria na busca da qualidade do ensino, Paro ( 2018 ) propõe uma perspectiva de participação da família junto à construção de gestão democrática de escola, pois assim geraria maiores benefícios para os sujeitos que constituem o espaço escolar e seu entorno.
O autor nos afirma que “ é possível imaginar um tipo de relação que não consista simplesmente de uma ‘ ajuda ’ gratuita dos pais à escola ” ( PARO, 2018, p. 39 ), mas a efetiva participação aos preceitos pedagógicos da instituição na busca de contribuir para a melhoria do ensino.
Pensar na relação democrática dentro do espaço escolar possibilita refletir que para a materialidade dessa relação é necessário compreender os problemas que envolvem esse vínculo família e escola, mensurar acerca dos valores que transitam nessa correlação e avaliar suas potencialidades.
Uma das principais características levantadas pelo autor ao tratar da relação família e escola é o \textbf{afeto} constituído junto aos seus alunos, sejam \textbf{crianças} ou adolescentes.
Esse \textbf{afeto} se constituiu como elemento importante no estabelecimento de referência educacional para a família e, principalmente, na constituição do respeito entre professor e aluno.
Respeito não apenas a sua condição de \textbf{criança}, que deve ser \textbf{cuidada}, protegida e tratada com carinho, mas também a seu direito de apropriar -se da cultura e de manifestar -se, sem constrangimentos deletérios, seu pensamento e sua emoção.
O \textbf{afeto} supõe empatia e compromisso do educador com o educando, com a preocupação de reforçar a condição de sujeito deste, estabelecendo uma relação humana que não seja fria e exterior, ocupada apenas em oferecer conhecimento para serem apreendidos, mas sim calorosa e cúmplice da própria formação de personalidade do educando ( PARO, 2007, p. 52 ).
O autor revela ainda que é importante compreender esse \textbf{afeto} não no sentido piegas, de “ autoajuda ” ao professor, mas entender o \textbf{afeto} de forma a auxiliá -lo no melhor desenvolvimento das suas atividades de ensino junto às suas \textbf{crianças} e/ou adolescentes.
Ao discorrer sobre a relação família e escola, esta necessita compreender e \textbf{acolher} a organização familiar que historicamente vem se constituindo, particularmente as do século XXI e, ainda considerar que as \textbf{crianças} são sujeitos históricos pertencentes a etnias e a povos diversos, de culturas diferenciadas.
Nesse sentido, ao propor maior interação entre família e a escola na busca de estabelecer uma relação mais forte e constante, essa relação deve estar pautada acima de tudo no respeito, na superação de preconceitos ou de estigmas evitando dissabor e dor às \textbf{crianças} que pertencem a diferentes famílias que hoje constituem a sociedade brasileira. [ 20 ] : O uso das aspas se dá pelo fato de a \textbf{criança} na verdade não deixa sua casa, mas a família a leva para o espaço da escola.

% matched lemmas: acolher, afeto, brincar, creche, criança, cuidado, demonstrar, desejo, infantil
\end{textsample}
